\section{问题一：确定天气下的时空网络优化}

\subsection{时间扩展网络 MILP}

对每一天 $t=1,\ldots,T$ 和有向边 $(i,j)\in E$ 定义二元变量 $x_{tij}$，表示当天由 $i$ 移动到 $j$；定义 $h_{ti}$ 表示在 $i$ 停留，$q_{ti}$ 表示在矿山 $i$ 挖矿，$y_{ti}$ 表示日末位于 $i$。每日恰执行一个行动：
\begin{equation}
\sum_{(i,j)\in E}x_{tij}+\sum_i h_{ti}+\sum_{i\in V_M}q_{ti}=u_{t-1},
\label{eq:one-action}
\end{equation}
其中 $u_t$ 表示第 $t$ 日尚未到达终点。位置流守恒为
\begin{equation}
y_{tj}=h_{tj}+q_{tj}+\sum_{i:(i,j)\in E}x_{tij},\qquad
\sum_j y_{tj}=u_t+\delta_t,
\label{eq:location-flow}
\end{equation}
其中 $\delta_t=1$ 表示当天首次到达终点。沙暴日直接令全部 $x_{tij}=0$，到达矿山当日不能挖矿由前一日位置与 $q_{ti}$ 的连接约束保证。

天气已知后，每日资源消耗是行动变量的线性组合。以水为例，
\begin{equation}
w_t=w_{t-1}+b^w_t-c^w_{z_t}
\left(2\sum_{(i,j)}x_{tij}+\sum_i h_{ti}+3\sum_{i\in V_M}q_{ti}\right).
\label{eq:milp-water}
\end{equation}
食物同理。购买变量只在起点第0日或村庄开放；起点仅允许一次购买，村庄价格乘2。现金平衡写为
\begin{equation}
m_t=m_{t-1}+B\sum_{i\in V_M}q_{ti}
-2p_w b^w_t-2p_f b^f_t,
\label{eq:cash}
\end{equation}
起点初购单价不乘2。负重约束\cref{eq:capacity}在初购和每次补给后均检查。目标最大化\cref{eq:terminal}。由于位置、行动、购买、库存和现金全部线性，模型是混合整数线性规划；其网络流骨架与线性优化理论一致\cite{bertsimas1997}。

\subsection{求解、回放与最优性证书}

模型由 HiGHS 1.15.1 求解。HiGHS 的学术引用依据其官方说明采用 Huangfu 与 Hall 的并行对偶修正单纯形论文\cite{huangfu2018}。求解器目标值返回后，不直接把变量表作为答案，而是将其转换成逐日计划，交由独立规则模拟器重新执行以下检查：邻接、沙暴移动、到矿当日挖矿、购买位置、采购现金、负重、资源非负和到达后停止。

\begin{table}[H]
\centering
\caption{第一、二关全局最优解摘要}
\label{tab:q1-summary}
\begin{tabularx}{0.92\textwidth}{C{1.3cm} C{2.0cm} C{2.0cm} C{1.6cm} C{2.2cm} Y}
\toprule
关卡 & 初始水/食物 & 最优终值/元 & 到达日 & 求解状态 & 路线与关键行为\\
\midrule
第一关 & 178/333 & \textbf{10470} & 24 & Optimal & 经15村庄、12矿山，挖矿7天，两次村庄补给\\
第二关 & 130/405 & \textbf{12730} & 30 & Optimal & 经39、62村庄和55矿山，累计挖矿12天\\
\bottomrule
\end{tabularx}
\end{table}

\subsection{第一关策略解释}

最优路径为
\[
1\to25\to26\to23\to21\to9\to15\to13\to12,
\]
在第11--16、19日于12区挖矿；随后经 $13\to15\to9\to21\to27$ 到达。第8日到15村庄后购买163箱水，第21日再次到15村庄购买36箱水和16箱食物。第一段采购主要修正水相对食物的消耗不平衡；第二段补给把终点剩余压到0，从而避免以全价或双倍价买入后只按半价回收。

矿山路线比最短路更长，但7天基础收益共7000元，足以覆盖绕行、挖矿三倍消耗和村庄溢价。求解器自动选择在沙暴第11日挖矿，却在连续沙暴第17、18日停留；原因不是沙暴不能挖矿，而是后两日继续挖矿会使返程资源和补给现金条件恶化。

\subsection{第二关策略解释与数据审计}

第二关先经 $1\to9\to18\to26\to35\to36\to37\to38\to39$，在39村庄分两日补水；再经47到55矿山。第14--18日挖矿，第19日到62村庄补给201箱水和86箱食物，第20日返回55矿山，第21--28日中继续挖矿7天，最后经56到64。

原公开参考轨迹把第15--18日水量依次记为110、86、56、26箱；按高温/沙暴挖矿三倍消耗应为100、76、46、16箱。第19日先完成高温移动，再购买201箱水和86箱食物，日末水量为201箱，与现金减少
\[
201\times2\times5+86\times2\times10=3730
\]
完全一致。修正后终值仍为12730元，且求解器找到不同路线但同目标值的另一最优解，说明本关存在多重最优。

\begin{figure}[H]
  \centering
  \includegraphics[width=0.84\textwidth]{known_weather_ledger.pdf}
  \caption{第一、二关最优策略的逐日库存、现金、天气与关键行动}
  \label{fig:q1-ledger}
\end{figure}

\begin{resultbox}{问题一结论 \evidence{E1}}
第一关最优终值10470元、第二关12730元，均由求解器全局最优状态和独立逐日回放共同支持。图\ref{fig:q1-ledger}同时证明补给、挖矿、现金与库存能够闭合，而非只有最终数值正确。
\end{resultbox}

\subsection{一般情况下的最优策略}

对任意地图、确定天气和有限截止日，只需替换邻接表、特殊节点与参数，即可复用上述时间扩展网络。若规模扩大，可先用到终点和矿山的最短距离做可达性剪枝，再对采购量建立更紧的上下界；若所有资源和现金保持整数，可直接使用分支定界给出最优间隙。故“一般策略”不是固定路线，而是：在完整天气路径下解联合位置—库存—现金 MILP，再用独立模拟器验证。
